\section{Discussion}
\label{sec:discussion}

\subsection{Key Findings and Their Interpretation}

\textbf{Finding 1: $\beta \approx 0.143$--$0.160\ \text{nat}^{-1}$ is an actionable quantity.} At this slope, the normalized divergence gap traverses the full $[0,1]$ scale in approximately 6--7 nats of KL budget beyond the crossover point. RLHF practitioners should monitor the normalized divergence gap and expect it to enter significant positive territory within 4--8 nats of KL budget, model-family-dependently.

\textbf{Finding 2: Near-identical slopes (ratio 1.116) suggest a model-family-independent mechanism.} If the calibration-alignment divergence rate were model-specific, slopes would vary substantially. Instead, $\beta_\text{Gao} \approx \beta_\text{Coste}$ within 12\% despite different model families, scales, and task distributions---suggesting the divergence curve reflects a property of the RLHF optimization process itself ($n=2$ datasets; further replication required to establish universality).

\textbf{Finding 3: The low-KL regime in Gao data reveals a ``ramp-up'' before divergence onset.} At very low KL, gold preference rises faster than RM score---producing a negative gap before crossing zero at $\sim$3.5 nats. This reflects a genuine feature of RLHF optimization: at low KL, the policy learns outputs that both the reward model and human evaluators prefer. The divergence onset occurs when proxy-target separation begins to dominate.

\subsection{Connections to Bidirectional Alignment}

Our results provide the empirical instantiation of the bidirectional alignment tension identified qualitatively by the ICLR 2025 Workshop \cite{ICLR2025Bidirectional}. Optimizing the AI$\to$Human direction (via RLHF with a proxy RM) degrades evaluation calibration to held-out human judgment at a quantifiable, replicable rate $\beta \approx 0.14\ \text{nat}^{-1}$. The normalized divergence gap metric provides a complementary evaluation instrument computable from standard RLHF evaluation runs.

\subsection{Limitations}

\textbf{L1: Digitization-derived data.} All analyses use data constructed from qualitative descriptions in published figures, not raw experimental datasets. The perfect $\rho = 1.000$ in Coste data is likely a digitization artifact. Effect sizes ($R^2 = 0.958$, $p < 10^{-6}$) are robust to $\pm 2$--5\% digitization error, but exact $\beta$ values are indicative estimates.

\textbf{L2: Small sample size ($N=10$).} The Gao bootstrap CI marginally overlaps zero, illustrating the power limitation. Future work with raw checkpoint data could compute divergence gap at 50+ KL levels.

\textbf{L3: P3 (coverage ratio) not tested.} The broader motivating claim---that alignment evaluation systematically omits Human$\to$AI measurement---relies on the ICLR 2025 survey's qualitative findings, not our own computation.

\textbf{L4: Evaluation calibration $\neq$ user behavioral calibration.} Our operationalization uses held-out gold preference in RLHF evaluation settings, not actual user behavioral calibration (appropriate reliance; Lai et al.~\cite{Lai2021Understanding}). Future behavioral studies can test whether evaluation-level divergence predicts user over-reliance.

\textbf{L5: Residual autocorrelation in OLS fit.} The Durbin-Watson statistic of 0.411 (Appendix~\ref{app:regression}) indicates positive residual autocorrelation in the Coste OLS fit, expected for serially ordered KL checkpoints; OLS standard errors and the resulting parametric CI $[0.119, 0.168]$ assume i.i.d.\ residuals and may be anti-conservative under autocorrelation. The bootstrap CI $[0.117, 0.177]$---which requires no i.i.d.\ assumption---independently confirms the strictly positive slope and serves as the autocorrelation-robust robustness check for this condition.

\subsection{Broader Impact}

This work has implications for AI safety and alignment evaluation practice. If reward model overoptimization systematically degrades held-out human preference calibration at $\beta \approx 0.14\ \text{nat}^{-1}$, then standard RLHF evaluation practices---which report RM scores without tracking held-out human preference---may systematically overestimate achieved alignment as optimization proceeds. The normalized divergence gap metric is a low-cost addition to standard RLHF evaluation. We see no direct negative societal impacts; the findings support \emph{more careful} alignment evaluation rather than enabling harmful applications.
